\section*{第 \ref{ch:g12:synthesis-strategy} 章——合成策略与绿色化学}

\begin{solution}{exo:g12:synthesis-strategy:1}
\omterm{def:g11:dissolution:solvation}{水合}：$46.0/(28.0 + 18.0) = \qty{100}{\percent}$。发酵：
$2 \times 46.0/180.0 = \qty{51.1}{\percent}$。
\end{solution}

\begin{solution}{exo:g12:synthesis-strategy:2}
两个基团，甘氨酸氨基上的 \ce{Z} 和丙氨酸羧基上的甲酯，
都在第 1 步引入、第 3 步除去。
\end{solution}

\begin{solution}{exo:g12:synthesis-strategy:3}
原则 5：使用更安全的\omterm{def:g6:solutions-and-solubility:solute}{溶剂}和反应条件。
\end{solution}

\begin{solution}{exo:g12:synthesis-strategy:4}
乙醇过量使 $Q < K$，反应更多地向正向进行。
\omterm{def:g12:catalysis:catalyst}{催化剂}对正、逆两个方向的反应同等程度地加速：
\omterm{def:g10:reaction-progress-table:system}{最终状态}相同，只是更快达到。
\end{solution}

\begin{solution}{exo:g12:synthesis-strategy:5}
可以：只有酮被还原为仲醇基团，
酯基保持不变。
\end{solution}

\begin{solution}{exo:g12:synthesis-strategy:6}
增加酯的产量：使一种\omterm{def:g7:chemical-reactions:reactant}{反应物}过量，或在生成水时将其移除（用水阱）。
加快反应：\omterm{def:g10:synthesis-yield:reflux}{加热回流}，加入酸性\omterm{def:g12:catalysis:catalyst}{催化剂}。
\end{solution}

\begin{solution}{exo:g12:synthesis-strategy:7}
用 \ce{NaBH4} 还原时：
\[ \ce{CH3-CH(OH)-CH2-CH2-CO-O-CH3} . \]
用 \ce{LiAlH4} 还原时：
\[ \ce{CH3-CH(OH)-CH2-CH2-CH2OH} \]
（戊-1,4-二醇），并由酯的另一部分生成甲醇。
\end{solution}

\begin{solution}{exo:g12:synthesis-strategy:8}
加成反应（\qty{100}{\percent}）、酯化反应
（\qty{83}{\percent}）和布洛芬的三步\omterm{def:g10:synthesis-yield:synthesis}{合成}路线
（\qty{77}{\percent}）。$77.4/40.0 = 1.9$：几乎是其两倍。
\end{solution}

\begin{solution}{exo:g12:synthesis-strategy:9}
$180.0/(138.0 + 102.0) = \qty{75.0}{\percent}$。乙酸：
$60.0/180.0 = \qty{0.333}{kg}$，每千克阿司匹林。
\end{solution}

\begin{solution}{exo:g12:synthesis-strategy:10}
加热加快了反应速率；冷凝器使蒸气回流到烧瓶中，
使\omterm{def:g7:chemical-reactions:reactant}{反应物}和\omterm{def:g6:solutions-and-solubility:solute}{溶剂}都不会损失。加热提高的是
反应速率，而不是最终\omterm{def:g10:synthesis-yield:yield}{产率}。
\end{solution}

\begin{solution}{exo:g12:synthesis-strategy:11}
六步对三步。原则：防止废物产生（1）、最大化
\omterm{def:g12:synthesis-strategy:atom-economy}{原子经济性}（2）、使用\omterm{def:g12:catalysis:catalyst}{催化剂}（9）；另外，步骤更少、能耗更低（6）。
\end{solution}

\begin{solution}{exo:g12:synthesis-strategy:12}
$0.80 \times 0.70 \times 0.90 = 0.504$：\qty{50.4}{\percent}。
起始原料：$0.10/0.504 = \qty{0.20}{mol}$。
\end{solution}

\begin{solution}{exo:g12:synthesis-strategy:13}
保护丙氨酸的氨基（\ce{Z}）和甘氨酸的羧基（甲酯）；
形成丙氨酸游离羧基与甘氨酸游离氨基之间的\omterm{def:g11:functional-groups:amine}{酰胺}键；
再除去两个\omterm{def:g12:synthesis-strategy:protecting-group}{保护基}。三个阶段，
若将每次保护和脱保护都计算在内，即五步反应。
\end{solution}

\begin{solution}{exo:g12:synthesis-strategy:14}
该\omterm{def:g12:catalysis:catalyst}{催化剂}符合原则 9。但 \qty{250}{\celsius} 不符合原则 6，
苯不符合原则 3 和原则 5，\omterm{def:g12:synthesis-strategy:atom-economy}{原子经济性}只有
\qty{35}{\percent}，不符合原则 2：$65/35 = \qty{1.9}{kg}$ 废物
每千克产物。一项绿色特征，
并不能使一个过程成为绿色过程。
\end{solution}

\begin{solution}{exo:g12:synthesis-strategy:15}
% ledger: crit:CO2-T, crit:CO2-p
\qty{40}{\celsius} 和 \qty{100}{bar} 均高于
二氧化碳的临界点（\qty{31}{\celsius}，\qty{73.8}{bar}）。压强是大气压强的一百倍，
因此容器必须是厚钢制的。二氧化碳无毒、不可燃，
且会变回气体、不留残渣，可以重复使用；
也不会有含氯\omterm{def:g6:solutions-and-solubility:solute}{溶剂}逸出。
\end{solution}

\begin{solution}{pb:g12:synthesis-strategy:1}
% ledger: ibu:bhc-ae-recovered
\textbf{1.} C：$10 + 4 + 1 = 15 = 13 + 2$；H：$14 + 6 + 2 = 22 = 18 + 4$；
O：$3 + 1 = 4 = 2 + 2$。

\textbf{2.} 乙酸，\ce{C2H4O2}。

\textbf{3.} 氮、氯和钠：布洛芬中只含有碳、
氢和氧。

\textbf{4.} 较新的那一种：原则 9，使用\omterm{def:g12:catalysis:catalyst}{催化剂}，
而非化学计量的\omterm{def:g7:identifying-substances:characteristic-test}{试剂}。

\textbf{5.} $13 \times 12.0 + 18 \times 1.0 + 2 \times 16.0 =
\qty{206.0}{g/mol}$。

\textbf{6.} $134.0 + 102.0 + 122.5 + 68.0 + 19.0 + 33.0 + 36.0 =
\qty{514.5}{g/mol}$。

\textbf{7.} $206.0/514.5 = \qty{40.0}{\percent}$.

\textbf{8.} $134.0 + 102.0 + 2.0 + 28.0 = \qty{266.0}{g/mol}$.

\textbf{9.} $206.0/266.0 = \qty{77.4}{\percent}$.

\textbf{10.} $(206.0 + 60.0)/266.0 = \qty{100}{\percent}$ 理论上，
实际中超过 \qty{99}{\percent}。

\textbf{11.} $514.5/206.0 = \qty{2.50}{t}$.

\textbf{12.} $2.50 - 1 = \qty{1.50}{t}$.

\textbf{13.} $266.0/206.0 = \qty{1.29}{t}$.

\textbf{14.} $60.0/206.0 = \qty{0.29}{t}$.

\textbf{15.} $1.50 - 0.29 = \qty{1.21}{t}$.

\textbf{16.} $0.90^6 = 0.53$: \qty{53}{\percent}.

\textbf{17.} $0.90^3 = 0.73$: \qty{73}{\percent}.

\textbf{18.} 每一步都要用到\omterm{def:g6:solutions-and-solubility:solute}{溶剂}，并进行分离和提纯，
也会损失部分产物：步骤越少，这类损失越少。

\textbf{19.} 几乎没有：其唯一的副产物得到了利用。

\textbf{20.} 三步路线每生产一吨布洛芬，可少产生约 \qty{1.5}{t} 废物，
（即便乙酸被当作废物丢弃，也仍可少产生约 \qty{1.2}{t}
）。
\end{solution}
